
\section{EEL 5544 : Motivation}
\begin{itemize}
  \item[*] What is random?
  \item[-] Some behavior that we counot perfectly predict (as oppose to deterministic)
  \item[*] Why randan?\\
Theologian : We can't know God's mind! ("nature"\\
Quantum physicist = the laus of quantum physics are fundamentally random.
  \item[*] What use of random?\\
Atthough cannot parfectly determine, still want to be able to obtain save degree of specifications.
  \item[*] Probability Theory is the theory that systewatically studies e characterizes randomness.
  \item[*] We will go through an engineering (operational, can less about wath vigor) trostment in this course.
  \item[*] Important concept (more like "terminology").\\
Randon experiment: An experiment whose atcomes are well randon!
\end{itemize}

\section{How should we study prob.?}
\begin{enumerate}
  \item Intuition\\
E.g.) Dr. Wong is a nice gry, the will probably give me a good grade in EEL 5544 .
\end{enumerate}

\begin{itemize}
  \item[-] Not systematic , quessing, cannot quantify.
\end{itemize}

\begin{enumerate}
  \setcounter{enumi}{1}
  \item Relative Frequency/Counting
\end{enumerate}

Lig.) Dr. Worg taught EEL5544 3 times in the past There were altogether 200 students who took Pr. War classes. Out of these 200 students, 50 of them got the grode "A". Thas the prob. of getting \(A\) in Dr. Wing's class is \(\frac{50}{200}=\frac{1}{4}\).

\begin{itemize}
  \item[-] Assume equally likely. this may not be the cate.
  \item[-] What about infinite possiblities.
\end{itemize}

\begin{enumerate}
  \setcounter{enumi}{2}
  \item Axiomatic throy
\end{enumerate}

\begin{itemize}
  \item[-] Identify a few basic axious that fit intertion
  \item[-] Build wath ansisteme system based on axioms
  \item[-] Utilize measure thery (a generalization of integration
  \item[-] Developed by Kolnogorou in conly 1900 s
\end{itemize}

